\documentclass[11pt]{article}

\usepackage[margin=1in]{geometry}
\usepackage[T1]{fontenc}
\usepackage{lmodern}
\usepackage{amsmath}
\usepackage{amssymb}
\usepackage{graphicx}
\usepackage{hyperref}

\hypersetup{colorlinks=true, linkcolor=blue, citecolor=blue, urlcolor=blue}

\newcommand{\CC}{\mathbb{C}}
\newcommand{\RR}{\mathbb{R}}
\newcommand{\Alg}{\mathcal{A}}
\newcommand{\EA}{E_{\mathcal{A}}}
\newcommand{\SA}{S_{\mathcal{A}}}
\newcommand{\rSM}{r_{\mathrm{SM}}}
\newcommand{\rY}{r_{Y}}

\title{Gauge-Invariant Algebraic Entropy in a Minimal Fermionic Toy:\\ A Single-Copy No-Go for Coupling Encoding}

\author{William T. Trevena\\[3pt]
{\small\itshape Independent Researcher (PhD, ISE, University of Florida)}\\
{\small\texttt{trevenaw7@gmail.com}}}

\date{June 2026}

\begin{document}

\maketitle

{\small\itshape\sloppy Acknowledgment: Portions of this manuscript were developed with the assistance of Claude (Anthropic), an AI language model, which contributed to literature synthesis, mathematical exposition, drafting, and the computational experiments. All numerical claims are asserted at runtime by the published code (\texttt{direction\_A/run\_direction\_A.py} and \texttt{direction\_A/strengthen.py}, deterministic seeds): the scripts fail if any verification breaks. Consistent with COPE authorship guidelines, the AI system is acknowledged as a tool rather than listed as a coauthor.\par}

\begin{abstract}
Whether an entanglement functional carries physical content in a gauged system turns on a question that can be settled by computation: is the functional invariant under the gauge action? This paper presents a general finite-dimensional method for deciding it --- compute the commutant of the unitary gauge action, identify its Wedderburn block structure, and evaluate the entropy relative to that gauge-invariant algebra, in the algebra-relative framework of Barnum--Knill--Ortiz--Viola and Zanardi, the discrete cousin of the Casini--Huerta center decomposition for gauge fields. The machinery applies unchanged to any finite-dimensional gauge action; we execute it end to end in a prominent three-qubit arena: the eight-dimensional Jordan--Wigner/Fock realization of the one-generation division-algebraic ideal $(\nu, \bar d, u, e^+)$, where a predecessor analysis \cite{companion} showed that gauge transformations of the fermionic modes are not local unitaries of the qubit factorization --- so single-qubit entropies are gauge-frame dependent and the question ``is the Standard Model coupling hierarchy encoded in three-qubit vacuum entanglement?'' was ill-posed as stated. The commutants are derived analytically --- the Fock space is the exterior algebra $\Lambda^0\CC^3 \oplus \Lambda^1\CC^3 \oplus \Lambda^2\CC^3 \oplus \Lambda^3\CC^3$, on which Schur's lemma fixes each commutant exactly (Theorem~1, proved in Appendix~B) --- and confirmed numerically: gauging $U(3)$ yields the abelian algebra $\CC^4$ of charge-sector projectors; $SU(3)$ yields $M_2 \oplus \CC \oplus \CC$, with the $M_2$ block on $\mathrm{span}\{\nu, e^+\}$; $U(1)$ yields $M_1 \oplus M_3 \oplus M_3 \oplus M_1$. Findings, each asserted at runtime by the published code (deterministic seeds): (i) the \emph{single-copy} gauge-invariant data of any pure state under $U(3)$ is its electric-charge-sector distribution alone; (ii) for every pure state and all three algebras the quantum piece of the algebraic entropy vanishes identically (Theorem~2) --- the toy's gauge-invariant entropy is 100\% classical center term, not distillable entanglement: the discrete analog of the edge-mode contribution in lattice gauge theory (``all edge, no bulk'' --- with the caveat, developed in Section~7.2, that the physical status of the edge term is itself contested \cite{chmp}); (iii) the predecessor's one surviving static ``match'' --- the weighted-W coupling-matching curve --- is pure gauge (Theorem~3): every state on it has $\SA = 0$ for all three algebras and is gauge-equivalent to a single Fock basis state, with $W = G(U)|001\rangle$ exhibited explicitly; (iv) the algebra-canonical rotor vacuum's invariant data is exactly $(1/2, 3/14, 3/14, 1/14)$ with $p_1 = p_2$, i.e.\ two independent invariant parameters --- fewer than three couplings; (v) an exhaustive search over 60 sector-assignment $\times$ feature-map combinations finds nothing within 0.05 of $\rSM = 1.8174$, the best unprincipled miss being $r = 2$ exactly; (vi) within this one-ideal, number-conserving three-mode realization, no $SU(2)_L$ action is present --- the mode gauge group is $U(3) \cong (SU(3)_c \times U(1)_{\mathrm{em}})/\mathbb{Z}_3$ --- so a two-ideal/chiral construction is required before the three-coupling question can even be posed. The scope of the no-go is \emph{single-copy}, and deliberately so: it closes the entropy version of the question and the single-copy invariant-data version --- everything measurable on one copy of the state under the charge superselection rule --- and leaves the multi-copy invariant ring open as the concrete next step (Section~8). The no-go concerns \emph{determination}, not \emph{reconstruction}: it rules out an algebra-dictated map from the kinematic invariant data of a single state to three couplings, while in two-dimensional Yang--Mills --- where dynamics supplies a known one-parameter family of flux-sector distributions --- the coupling is exactly reconstructible from precisely this kind of center data \cite{ym2companion}. For any version of the question to survive, two things would have to exist: a two-ideal construction on which $SU(2)_L$ acts --- specified in Section~9, and since carried out in a drafted companion manuscript \cite{twoideal} --- and a principle linking sector statistics to couplings, which remains absent.
\end{abstract}

\section{Introduction}

Any proposal that reads physics out of an entanglement functional in a gauged system faces a prior question: is the functional gauge-invariant at all? In finite dimensions that question can be settled mechanically, and this paper is, before anything else, a demonstration of the method: given a unitary gauge action on a finite-dimensional Hilbert space, compute the commutant --- the algebra of gauge-invariant observables --- identify its Wedderburn block structure, and evaluate entropy relative to that algebra. The resulting functional is gauge-invariant by construction; the functional it replaces can be tested for frame dependence directly; and the block structure of the commutant tells one, in advance, how much of the invariant entropy can ever be quantum rather than classical. None of this machinery is specific to the system studied here --- it applies unchanged to any finite-dimensional gauge action. What this paper contributes is the method executed exactly, end to end, on a prominent target, together with the no-go results that the execution delivers.

The target arises as follows. A predecessor paper \cite{companion} asked whether the Standard Model gauge coupling hierarchy is encoded in the entanglement entropies of a three-qubit vacuum state, motivated by the convergence of the emergent-spacetime and division-algebraic programs through the Szangolies construction \cite{szangolies}. Its answer was a battery of no-go results, capped by a diagnosis: the question is \emph{ill-posed}. Recognizing Furey's $\mathrm{Cl}(6)$ construction \cite{furey14, furey18, stoica} as the three-qubit operator algebra under the Jordan--Wigner transformation \cite{jordanwigner} makes electric charge the Hamming-weight grading of the computational basis --- and exposes the obstruction that gauge transformations act on the fermionic \emph{modes}, not on the qubit tensor factors. A color discrete Fourier transform maps a separable single-quark state, with single-qubit entropies $(0, 0, 0)$, to the W state, with entropies $(0.918, 0.918, 0.918)$. The program's central quantity was frame decoration.

That paper closed by identifying the construction of a gauge-invariant entanglement functional as the central open problem (\cite{companion}, Secs.~5.3 and~8). This paper completes that program for the minimal toy, and reports what the completed version says.

The functional is not new and we do not claim it as such: entanglement relative to a distinguished $*$-subalgebra of observables is the ``generalized entanglement'' of Barnum, Knill, Ortiz, and Viola \cite{barnum}, the algebra-relative subsystem notion of Zanardi \cite{zanardi}, and --- in the form most relevant here --- the algebraic definition that Casini, Huerta, and Rosabal introduced for gauge fields \cite{chr}, whose classical center term Donnelly identified with lattice-gauge edge modes \cite{donnelly12, donnelly14}. Our contribution is to \emph{execute} this definition exactly, end to end, in the one system where the question of \cite{companion} lives: the eight-dimensional one-generation ideal $(\nu, \bar d, u, e^+)$ in its Jordan--Wigner three-qubit presentation, under the three natural gauge choices $U(3)$, $SU(3)$, and $U(1)$. Everything is finite-dimensional, so there are no regularization choices to argue about; the commutants are computed numerically and certified, the entropies have closed forms, and every claim in this paper is asserted at runtime by the published code.

\subsection{Contributions}

\begin{enumerate}

\item \textbf{The gauge-invariant algebras, proved and verified numerically} (Theorem~1, Section~4; analytic proof by exterior-algebra decomposition and Schur's lemma in Appendix~B; independent numerical certification in Section~4 and Lemma~C.2). Gauging $U(3)$ leaves only the abelian algebra $\CC^4$ of charge-sector projectors; $SU(3)$ leaves $M_2 \oplus \CC \oplus \CC$ (dimension 6), the $M_2$ block living on $\mathrm{span}\{\nu, e^+\}$ --- the $\nu$--$e^+$ coherence is gauge-invariant when only color is gauged; $U(1)$ leaves $M_1 \oplus M_3 \oplus M_3 \oplus M_1$ (dimension 20). Under the full $U(3)$, the single-copy gauge-invariant data of any pure state is its charge-sector distribution and nothing else.

\item \textbf{All edge, no bulk} (Theorem~2, Section~5). For every pure state and all three algebras, the quantum piece of the algebraic entropy vanishes identically --- a two-line consequence of the block structure, verified numerically to $1.8 \times 10^{-15}$. The toy's entire gauge-invariant entropy is the classical center term, the discrete analog of the edge-mode contribution in lattice gauge theory. As a corollary we derive $S_{U(1)} = S_{U(3)}$ for every pure state, despite the $U(1)$-invariant algebra being five times larger.

\item \textbf{The headline: the predecessor's surviving ``match'' is pure gauge} (Theorem~3, Section~5.4). Every point of the weighted-W coupling-matching curve --- the one static construction in \cite{companion} that reproduced $\rSM = 1.8174$ --- has $\SA = 0$ for all three invariant algebras and is gauge-equivalent to a single Fock basis state: one quark in a rotated color frame. We exhibit $U \in SU(3)$ with $G(U)|001\rangle = |W\rangle$ exactly. The static coupling-matching analysis of \cite{companion}, Secs.~2--4, was matching frame artifacts.

\item \textbf{The well-posed question, answered} (Section~6). The gauge-invariant reformulation of the coupling-entropy question is stated and answered in the negative for four independent reasons, each checked in the accompanying code: the invariant sectors are charge sectors, not gauge sectors; the algebra-canonical rotor vacuum carries exactly two independent invariant parameters, $p = (1/2, 3/14, 3/14, 1/14)$ with $p_1 = p_2$; exhaustive search over sector-to-coupling assignments misses $\rSM$ (best unprincipled miss: $r = 2$ exactly, flagged as numerology bait); and, within this one-ideal, number-conserving three-mode realization, no $SU(2)_L$ action is present --- a deliberately scoped lemma, verified in the accompanying code (Lemma~4; Appendix~C.1), not a claim about Furey's framework in general --- so a three-coupling question cannot even be posed here.

\item \textbf{The bridge to the edge-mode literature, stated with its contested status on the surface} (Section~7.2). The center term computed throughout is the same object, in an 8-dimensional toy, as the classical flux-sector term of lattice gauge theory \cite{chr, donnelly12, donnelly14} and the edge-mode reading of Kabat's contact term \cite{kabat, dw15, dw16} --- and we record honestly, following Casini--Huerta--Mag\'an--Pontello \cite{chmp}, that the physical status of that term is itself contested in the continuum.

\item \textbf{What entropy cannot see, and what comes next} (Sections~8 and~9). The entropy is a lossy summary even of the single-copy invariant data: under $U(1)$ the invariant algebra is five times larger than under $U(3)$, yet the two entropies coincide for every pure state (Corollary~2.1). We therefore separate the closed entropy version of the question from the open invariant-data version, pose the multi-copy invariant ring (invariant theory, Hilbert series, Casimir cumulants) as the concrete next step, and specify the two-ideal extension --- the smallest arena on which $SU(2)_L$ acts --- as a step-by-step program; that computation has since been completed and is reported in a drafted companion manuscript \cite{twoideal} (Section~9).

\end{enumerate}

\subsection{Scope}

This paper is about an eight-dimensional toy. Nothing here bears on continuum gauge theory except by structural analogy, and the analogy's load-bearing limit is stated where it is used. The no-go results are exact within the toy and claim nothing beyond it; their value, if any, is that they close the minimal case cleanly and locate exactly what a non-minimal case would need.

\subsection{Three Entropies, Distinguished}

Three distinct entropy functionals appear in this paper and its predecessors, and every claim below is about exactly one of them. The table fixes the terminology once.

\begin{center}
\small
\begin{tabular}{p{3.4cm}p{5.3cm}p{5.3cm}}
\hline
object & definition & role in this paper \\
\hline
\textbf{qubit subsystem entropy} $S(\rho_i)$ & von Neumann entropy of a single-qubit reduced density matrix in the Jordan--Wigner qubit factorization & gauge-\emph{variant}; the quantity of \cite{companion}, shown only for contrast (Section~5.2, rightmost column) and revealed as frame artifact (Theorem~3) \\
\textbf{algebraic (algebra-relative) entropy} $\SA$ & Eq.~(\ref{eq:salg}): entropy relative to the gauge-invariant algebra $\Alg$ --- classical center piece $H(\{p_k\})$ plus quantum piece $\sum_k p_k S(\rho_k)$ & gauge-invariant by construction; the central object of Sections~3--8 \\
\textbf{flux-sector (edge) entropy} $H(\{p_k\})$ & Shannon entropy of the center / superselection-sector distribution --- the classical piece of $\SA$; in lattice gauge theory, the boundary-flux (edge-mode) term & equals \emph{all} of $\SA$ for pure states in this toy (Theorem~2); the bridge object of Section~7.2 and of the two-dimensional Yang--Mills analysis \cite{ym2companion} \\
\hline
\end{tabular}
\end{center}

\section{The Toy and Its Gauge Action}

\textbf{Status: construction inherited from \cite{companion}, Sec.~5; verified numerically here to stated residuals.}

The arena is $\CC^8 = \CC^2 \otimes \CC^2 \otimes \CC^2$ with three Jordan--Wigner fermionic modes
\begin{equation*}
f_1 = \sigma^- \otimes I \otimes I, \qquad
f_2 = Z \otimes \sigma^- \otimes I, \qquad
f_3 = Z \otimes Z \otimes \sigma^-,
\end{equation*}
satisfying the canonical anticommutation relations exactly (residual $0.0$ in our run). The computational basis is the Fock basis, graded by Hamming weight, and the weight is three times the electric charge on Furey's one-generation ideal \cite{furey14, furey18, stoica}: $\nu$ at weight 0, the $\bar d$ color triplet at weight 1, the $u$ triplet at weight 2, $e^+$ at weight 3, with $Q = N/3$ for $N$ the number operator.

A mode rotation $U \in U(3)$, $f_b^\dagger \mapsto \sum_a U_{ab} f_a^\dagger$, induces the number-conserving unitary $G(U)$ on $\CC^8$ acting blockwise on the weight sectors:
\begin{equation*}
\text{weight 0: } 1, \qquad
\text{weight 1: } U, \qquad
\text{weight 2: } \Lambda^2 U, \qquad
\text{weight 3: } \det U,
\end{equation*}
with $G(e^{i\theta}\mathbf{1}) = e^{i\theta N}$ on the $U(1)$ subgroup. The construction was verified to machine precision: $G$ unitary (residual $2.4 \times 10^{-15}$), homomorphism $G(U)G(V) = G(UV)$ ($6.8 \times 10^{-16}$), and the mode-transformation property $G(U) f_b^\dagger G(U)^\dagger = \sum_a U_{ab} f_a^\dagger$ ($1.9 \times 10^{-15}$).

Three gauge choices are considered, each physically motivated: the full $\boldsymbol{U(3)}$ (color and electromagnetism together --- the largest group of number-conserving mode rotations); $\boldsymbol{SU(3)}$ alone (gauge color, leave charge ungauged); and $\boldsymbol{U(1)}$ alone (charge superselection without color gauging). The group $U(3) \cong (SU(3)_c \times U(1)_{\mathrm{em}})/\mathbb{Z}_3$ is the toy's entire mode gauge group; Section~6.5 explains why no $SU(2)_L$ action is present in this realization --- and what construction would be required before one could be.

\section{Entropy Relative to a \texorpdfstring{$*$}{*}-Algebra}

\subsection{Definition}

Let $\Alg$ be a unital $*$-subalgebra of $M_d(\CC)$. By the Wedderburn--Artin theorem there is a unitary change of basis under which
\begin{equation*}
\mathcal{H} \cong \bigoplus_k \CC^{n_k} \otimes \CC^{m_k}, \qquad
\Alg = \bigoplus_k M_{n_k}(\CC) \otimes \mathbf{1}_{m_k},
\end{equation*}
with commutant $\Alg' = \bigoplus_k \mathbf{1}_{n_k} \otimes M_{m_k}$ and center $Z(\Alg) = \Alg \cap \Alg' \cong \CC^{\#\mathrm{blocks}}$, spanned by the minimal central projections $P_k$. The reduced state of $|\psi\rangle$ on $\Alg$ is the conditional expectation $\EA(|\psi\rangle\langle\psi|)$ --- equivalently, the unique state $\omega$ on $\Alg$ with $\omega(a) = \langle\psi|a|\psi\rangle$. Writing
\begin{equation*}
p_k = \langle\psi|P_k|\psi\rangle, \qquad
\rho_k = \frac{1}{p_k} \mathrm{Tr}_{m_k}\!\left(P_k |\psi\rangle\langle\psi| P_k\right),
\end{equation*}
the \textbf{algebraic entropy} (algebra-relative entropy) is
\begin{equation}
\SA(\psi) = \underbrace{H(\{p_k\})\vphantom{\sum_k}}_{\text{classical (center) piece}} \;+\; \underbrace{\sum_k p_k S(\rho_k)}_{\text{quantum piece}},
\label{eq:salg}
\end{equation}
in bits throughout. This is the generalized entanglement of \cite{barnum}, the algebra-relative notion of \cite{zanardi}, and --- crucially --- the discrete cousin of the Casini--Huerta--Rosabal definition for gauge fields \cite{chr}: when the algebra of a region in a lattice gauge theory has a nontrivial center (boundary electric fluxes), the entropy acquires exactly the $H(\{p_k\})$ term --- the Shannon entropy of the boundary-flux superselection sectors, nowadays read as the edge-mode contribution \cite{donnelly12, donnelly14}. The center term computed throughout this paper is the same object in an 8-dimensional toy.

\textbf{Terminology.} We say \emph{algebraic entropy} (equivalently: center entropy, algebra-relative entropy) rather than ``entanglement'' wherever the center term is the content. The reason is visible in Eq.~(\ref{eq:salg}) itself: $H(\{p_k\})$ is the Shannon entropy of superselection-sector weights --- classical data, not distillable entanglement \cite{donnelly14, chr, soni}. The names under which the general functional was introduced (``generalized entanglement'' \cite{barnum}, entanglement relative to an algebra \cite{zanardi}) are retained when citing those frameworks, but no claim of operational, distillable entanglement is intended or needed anywhere in this paper; the contested physical status of the center term is taken up in Section~7.2.

Two consistency limits: if $\Alg = B(\mathcal{H}_A) \otimes \mathbf{1}_B$ is a factor, $\SA = S(\rho_A)$ and the center piece is zero --- the standard entanglement entropy is recovered; if $\Alg$ is abelian (all $n_k = 1$), $\SA = H(\{p_k\})$ is purely classical.

\textbf{Pitfall (cross-checked numerically on every call).} $\SA$ is \emph{not} the von Neumann entropy of the $d \times d$ matrix $\EA(\rho)$; the two differ by $\sum_k p_k \log_2 m_k$, because $\EA(\rho)$ spreads each block uniformly over its multiplicity factor. The library computes $\SA$ both ways and asserts agreement to $10^{-8}$.

\subsection{Why This Is the Right Gauge-Invariant Functional}

If a compact group acts on $\mathcal{H}$ by unitaries $G(g)$, the gauge-invariant observables are exactly the commutant $\Alg = \{G(g)\}'$. Conjugation by any $G(g)$ fixes $\Alg$ pointwise, so $\EA(G\rho G^\dagger) = \EA(\rho)$: \textbf{$\SA$ is gauge-invariant by construction.} The quantity that was frame-dependent in \cite{companion} --- single-qubit entropy --- is replaced by one that provably cannot be. (Verified numerically below to $\sim$$10^{-15}$, alongside the $O(1)$ shifts of the single-qubit entropies under the same transformations.)

\textbf{Honest scope note.} $\EA(\rho)$ captures every statistic obtainable by \emph{single-copy measurements of gauge-invariant observables} --- under a charge superselection rule this is all that is operationally accessible \cite{brs}. The polynomial invariant ring of the state under $G$ is strictly larger: it contains multi-copy invariants (computed for the rotor vacuum in Section~8) that are constants of the $G$-orbit but not functions of $\EA(\rho)$. The well-posed question of Section~6 is phrased in terms of the operational single-copy data; the multi-copy ring is flagged as the remaining unexplored invariant content, not silently absorbed.

\section{The Gauge-Invariant Algebras}

\textbf{Status: proved (Theorem~1; Appendix~B) and independently computed numerically (seed 20260609; \texttt{direction\_A/results.json}).}

\medskip

\noindent\textbf{Theorem 1 (commutant structure).} \emph{Let $G$ be the mode-rotation representation of Section~2, acting on the Hamming-weight sectors of $\CC^8$ as $(1, U, \Lambda^2 U, \det U)$, and let $P_k$ denote the projector onto weight $k$. The commutants of the three gauge actions are exactly:}
\begin{align*}
\{G(U) : U \in U(3)\}' &= \mathrm{span}\{P_0, P_1, P_2, P_3\} \cong \CC^4 && \text{\emph{(abelian; dimension 4),}}\\
\{G(U) : U \in SU(3)\}' &\cong M_2 \oplus \CC \oplus \CC && \text{\emph{(dimension 6),}}\\
\{G(e^{i\theta}\mathbf{1}) : \theta \in \RR\}' &= M_1 \oplus M_3 \oplus M_3 \oplus M_1 && \text{\emph{(dimension 20),}}
\end{align*}
\emph{where the $M_2$ block in the second line is the full matrix algebra on $\mathrm{span}\{|000\rangle, |111\rangle\}$ and the third line is the algebra of all weight-preserving operators.}

\textbf{Proof.} Appendix~B: identify the Fock space with the exterior algebra $\Lambda^0\CC^3 \oplus \Lambda^1\CC^3 \oplus \Lambda^2\CC^3 \oplus \Lambda^3\CC^3$, on which $G(U)$ acts as $1 \oplus U \oplus \Lambda^2 U \oplus \det U$, and apply Schur's lemma to each gauge choice. \hfill$\square$

\medskip

\textbf{Numerical confirmation (independent of the proof).} For each gauge choice, 50 Haar samples $G(U_s)$ were drawn and the commutant $\{X : [X, G(U_s)] = 0\ \forall s\}$ computed as the SVD nullspace of the stacked maps $X \mapsto (\mathbf{1} \otimes G - G^{\mathsf{T}} \otimes \mathbf{1})\,\mathrm{vec}(X)$. Each basis was certified against 100 fresh samples (maximum commutator residual $< 1.2 \times 10^{-15}$), checked to be a unital $*$-algebra (closure residual $< 1.3 \times 10^{-15}$), and its Wedderburn structure read off from the center and per-block ranks --- assuming nothing from Theorem~1. The extracted structures agree with the theorem in every detail: dimensions, block supports, and centers. (Lemma~C.2 closes the loop from a third, independent side: for each gauge choice, exhibited exactly-commuting elements span a space whose dimension equals the numerically computed commutant dimension, so the identifications are exact even without Appendix~B.)

\begin{center}
\small
\begin{tabular}{p{1.5cm}p{0.9cm}p{6.1cm}p{1.0cm}p{4.2cm}}
\hline
gauged group & $\dim \Alg$ & structure (block [weight support]) & center & invariant data of a pure state \\
\hline
$\boldsymbol{U(3)}$ & \textbf{4} & $M_1[w0] \oplus M_1{\otimes}\mathbf{1}_3[w1] \oplus M_1{\otimes}\mathbf{1}_3[w2] \oplus M_1[w3]$ --- abelian $\CC^4$ & $\CC^4$ & $(p_0, p_1, p_2, p_3)$: \textbf{3 real parameters} \\
$\boldsymbol{SU(3)}$ & \textbf{6} & $\boldsymbol{M_2}{\otimes}\mathbf{1}_1[w0 \cup w3] \oplus M_1{\otimes}\mathbf{1}_3[w1] \oplus M_1{\otimes}\mathbf{1}_3[w2]$ & $\CC^3$ & $(p_1, p_2)$ + pure qubit in the $M_2$ block: \textbf{4 parameters} \\
$\boldsymbol{U(1)}$ & \textbf{20} & $M_1[w0] \oplus \boldsymbol{M_3}[w1] \oplus \boldsymbol{M_3}[w2] \oplus M_1[w3]$ & $\CC^4$ & 3 probabilities + two projective pure states in $\CC^3$: \textbf{11 parameters} \\
\hline
\end{tabular}
\end{center}

Readings, in turn:

\begin{itemize}

\item \textbf{$U(3)$: the commutant is abelian.} The four weight sectors carry the inequivalent irreps $(1, U, \Lambda^2 U, \det U)$, each with multiplicity one, so the invariant algebra is the sector projectors and nothing more. \emph{The only single-copy gauge-invariant data of any state is its charge-sector distribution $(p_0, p_1, p_2, p_3)$.} This is the precise, gauge-invariant version of \cite{companion}'s ``index ambiguity'' (Sec.~2.1 Remark there): there is no gauge-covariant tensor factorization because there are no gauge-invariant subsystems at all --- only charge sectors.

\item \textbf{$SU(3)$: the $\nu$--$e^+$ coherence survives.} Weights 0 and 3 are both $SU(3)$-trivial, so the commutant contains a full $M_2$ block on $\mathrm{span}\{|000\rangle, |111\rangle\} = \mathrm{span}\{\nu, e^+\}$ (numerical check: the $M_2$ central projection equals the projector onto indices $\{0, 7\}$; assertion threshold $10^{-9}$, measured residual $\leq 3.3 \times 10^{-15}$). When only color is gauged, the GHZ-type $\nu$--$e^+$ coherence is gauge-invariant data. Weights 1 and 2 (the $\mathbf{3}$ and $\bar{\mathbf{3}}$) remain inequivalent and contribute $\CC \oplus \CC$.

\item \textbf{$U(1)$: everything weight-preserving.} $G(e^{i\theta}\mathbf{1}) = e^{i\theta N}$, so the commutant is $M_1 \oplus M_3 \oplus M_3 \oplus M_1$ (dimension 20) --- the richest of the three, retaining the full within-sector color states. Parameter-count sanity check: a pure state in $\CC^8$ has 14 real parameters; gauging $U(1)$ erases exactly the 3 relative phases between the four sectors: $14 - 3 = 11$. \checkmark

\end{itemize}

So under $U(3)$ the invariant data has exactly 3 real parameters --- naive room for one per coupling. Section~6 is about why that room cannot be used.

\section{Gauge-Invariant Algebraic Entropy of the Named States}

\textbf{Status: computed; closed forms checked numerically to $10^{-12}$.}

\subsection{The Quantum Piece Vanishes Identically}

\textbf{Theorem 2 (all edge, no bulk).} \emph{For every pure state $|\psi\rangle \in \CC^8$ and each of the three algebras of Theorem~1, the quantum piece of Eq.~(\ref{eq:salg}) vanishes: $\SA(\psi) = H(\{p_k\})$.}

\textbf{Proof.} Every Wedderburn block of all three algebras has $n_k = 1$ or $m_k = 1$. If $n_k = 1$, then $\rho_k$ is a $1 \times 1$ state and $S(\rho_k) = 0$ trivially. If $m_k = 1$, the multiplicity factor is trivial, so $P_k|\psi\rangle$ is a vector in the block and $\rho_k = P_k|\psi\rangle\langle\psi|P_k / p_k$ is a rank-one projector: $S(\rho_k) = 0$. \hfill$\square$

(Numerical confirmation: maximum $|$quantum piece$|$ $= 1.8 \times 10^{-15}$ over the named states and a 200-state subsample of the $10^4$-state Haar ensemble run through the full block machinery; the remaining ensemble states were evaluated with a sector formula that agrees with the full machinery to $2.0 \times 10^{-15}$ on that subsample.)

We do not present Theorem~2 as a deep surprise: given the Wedderburn block structure of Theorem~1, the vanishing of the quantum piece for pure states follows directly, as the two-line proof shows; the theorem's value lies in its interpretation and in the bridge it provides to the gauge-theory edge-term literature (Section~7.2), not in the calculation itself.

In Casini--Huerta language: for pure states, this toy is \emph{all edge-mode term, no bulk term}. Any claim that ``vacuum entanglement encodes couplings'' in this framework is therefore, gauge-invariantly, a claim about classical charge statistics. Section~7.2 develops this bridge; here we note the scope limit on the surface: Theorem~2 is a statement about \emph{pure} states. For mixed states the $SU(3)$ and $U(1)$ algebras can acquire nonzero quantum pieces (the $U(3)$ entropy stays classical for all states because the algebra is abelian); and the first system in this program where a nonzero quantum piece could appear for pure states needs blocks with both $n_k, m_k > 1$ --- e.g.\ regional subalgebras, or the two-ideal extension of Section~9.

\textbf{Corollary 2.1 ($S_{U(1)} = S_{U(3)}$ for every pure state).} \emph{The $U(1)$ and $U(3)$ commutants have the same center --- the four weight projectors. By Theorem~2, both entropies equal $H(p_0, p_1, p_2, p_3)$.} \hfill$\square$

This holds even though the $U(1)$-invariant algebra is five times larger: gauging the extra $SU(3)$ changes \emph{which states are distinguishable} (the $U(1)$ data distinguishes W from $|001\rangle$; the $U(3)$ data does not) but not the entropy. Entropy is a very lossy summary of $\EA(\rho)$ --- a deflation worth keeping in view whenever a single number is asked to carry structural weight. Section~8 elevates this from caveat to result.

\subsection{The Table}

All entropies in bits; sector probabilities exact rationals (checked against the numerics to $10^{-10}$). The rightmost column shows the gauge-\emph{variant} single-qubit entropies for contrast.

\begin{center}
\small
\begin{tabular}{p{2.9cm}p{2.5cm}cccp{2.5cm}}
\hline
state & $(p_0, p_1, p_2, p_3)$ & $S_{U(3)}$ & $S_{SU(3)}$ & $S_{U(1)}$ & single-qubit $S$ (gauge-variant) \\
\hline
\textbf{rotor vacuum} $(\sqrt{7}\,e_0 - iu)/\sqrt{14}$ & $\mathbf{(1/2, 3/14,}$ $\mathbf{3/14, 1/14)}$ exact & \textbf{1.724408} & \textbf{1.413800} & 1.724408 & $(0.5917,$ $0.5917,$ $0.5917)$ \\
\textbf{W} $=$ color-uniform $\bar d$ & $(0, 1, 0, 0)$ & \textbf{0} & \textbf{0} & \textbf{0} & $(0.9183,$ $0.9183,$ $0.9183)$ \\
\textbf{GHZ} $= (\nu + e^+)/\sqrt{2}$ & $(1/2, 0, 0, 1/2)$ & \textbf{1} & \textbf{0} & 1 & $(1, 1, 1)$ \\
\textbf{weighted-W matching state} (\cite{companion}, Fig.~1c curve) & $(0, 1, 0, 0)$ & \textbf{0} & \textbf{0} & \textbf{0} & $(0.9995,$ $0.8580,$ $0.7802)$ \\
Haar ensemble ($10^4$), mean $\pm$ std & --- & $1.580 \pm 0.239$ & $1.398 \pm 0.162$ & $= S_{U(3)}$ & --- \\
\hline
\end{tabular}
\end{center}

{\sloppy \emph{Cross-paper note (weighted-W row).} The representative matching state quoted here, $S = (0.9995, 0.8580, 0.7802)$, is the first stored point of the exp2 matching curve (\texttt{weighted\_W.examples[0]} in \texttt{exp2\_results.json}, curve grid index 0). It is the same representative quoted in \cite{companion} when verifying its closed-form Lemma~2 (Sec.~4.1 there, the $g = 0.0779$ point). The separately stored curve midpoint, $S = (0.978, 0.895, 0.849)$ (\texttt{weighted\_W\_example}, grid index 21 of 43), is covered by the same verification: the two are different points on the same one-parameter $r_S = 1.8174$ matching curve, not a discrepancy --- both satisfy $r_S = \rSM$ to solver precision, both are verified weight-1-supported, and both are pure gauge, with $\SA = 0$ for all three invariant algebras (verified in \texttt{direction\_A/strengthen.py}).\par}

Exact closed forms for the rotor vacuum (checked numerically to $10^{-12}$):
\begin{align*}
S_{U(3)} = S_{U(1)} &= \tfrac{1}{2} + \tfrac{1}{2}\log_2 14 - \tfrac{3}{7}\log_2 3 = 1.724407817863 \text{ bits}, \\
S_{SU(3)} &= \tfrac{4}{7}\left(\log_2 7 - 2\right) + \tfrac{3}{7}\log_2\!\frac{14}{3} = 1.413799564606 \text{ bits},
\end{align*}
and the $SU(3)$-invariant qubit in the $M_2$ ($\nu/e^+$) block is pure, with block probability $4/7$ and state $\rho = \frac{1}{8}\begin{pmatrix} 7 & i\sqrt{7} \\ -i\sqrt{7} & 1 \end{pmatrix}$ --- the gauge-invariant $\nu$--$e^+$ coherence.

\textbf{Gauge invariance verified (the whole point).} For every named state, every algebra, and 20 random elements of the corresponding gauge group: $\max |\Delta \SA| < 3 \times 10^{-15}$ (assertion threshold $10^{-9}$), while the single-qubit entropies of the \emph{same} transformed states shift by up to 0.9226 bits. Bonus, also verified to $\sim$$10^{-15}$: $S_{SU(3)}$ and $S_{U(1)}$ are invariant under the \emph{full} $U(3)$ --- $G(U)$ normalizes the weight decomposition, and the extra $U(1)$ phases are isospectral on the $M_2$ block.

\subsection{State by State}

\textbf{The W state's qubit entanglement is pure gauge.} We exhibit $U \in SU(3)$ (third column $(1,1,1)/\sqrt{3}$) with $G(U)|001\rangle = |W\rangle$ exactly (construction residual $0.0$), and verify $\EA(W) = \EA(|001\rangle)$ for the $U(3)$ algebra (residual $2.6 \times 10^{-16}$). The paradox that opened the well-posedness diagnosis in \cite{companion}, Sec.~5.3 --- a separable quark state mapped to the ``entangled'' W by a color DFT --- is here resolved rather than flagged: the algebraic entropy assigns \textbf{zero to both}, because a single quark in a color-rotated frame is still a single quark. The 0.918 bits the W state carries in the qubit factorization is entirely frame.

\textbf{GHZ is one classical bit --- or nothing.} Under $U(3)$, which contains the electromagnetic $U(1)$, the GHZ state's invariant content is the charge distribution $(1/2, 0, 0, 1/2)$: exactly one bit of classical charge uncertainty, no quantum piece. Under $SU(3)$ only, the $\nu$--$e^+$ coherence sits inside the $M_2$ invariant block as a \emph{pure} qubit, and $S_{SU(3)}(\mathrm{GHZ}) = 0$. The ``maximal tripartite entanglement'' of GHZ is, gauge-invariantly, either one classical bit (charge superselection imposed) or no entropy at all (color-only gauging). Nothing about it is three-partite in any invariant sense. The toy cleanly separates color gauging from charge superselection here: the $\nu$--$e^+$ coherence is physical data if and only if the abelian factor is not gauged.

\textbf{The rotor vacuum} --- the unique ground state of the algebra-canonical Hamiltonian $H = i(L_u + 2R_u)$ of \cite{companion}, Appendix~A --- is the only state in the list with nontrivial invariant data, and that data is two numbers (Section~6.3). Its invariant entropy, 1.724 bits, is comfortably \emph{generic}: the Haar mean is $1.580 \pm 0.239$. Gauge-invariantly, the canonical vacuum is not even unusually entangled.

\subsection{The Headline: the Coupling-Matching Curve Collapses to Pure Gauge}

The only surviving static ``match'' of \cite{companion} was the weighted-W family: a continuous curve of states $a|001\rangle + b|010\rangle + c|100\rangle$ reproducing the Standard Model log-inverse-coupling gap ratio $\rSM = 1.8174$ exactly under the two-parameter logarithmic map (\cite{companion}, Eq.~1; Secs.~2.2--2.3, 4.1, Fig.~1c). Every state on that curve is a weight-1 state, so the following theorem annihilates the entire matching mechanism at once --- and covers the W state of Section~5.3 as the special case $a = b = c$.

\medskip

\noindent\textbf{Theorem 3 (the matching curve is pure gauge).} \emph{Every unit-norm weight-1 state $\psi = a|001\rangle + b|010\rangle + c|100\rangle$ --- in particular every point of the weighted-W matching curve, and the W state --- satisfies}

\emph{(i) $\psi = G(U)|001\rangle$ for some $U \in SU(3)$: the whole curve lies in a single gauge orbit, that of the Fock basis state $|001\rangle$;}

\emph{(ii) $\SA(\psi) = 0$ for all three invariant algebras of Theorem~1.}

\textbf{Proof.} (i) $SU(3)$ acts transitively on the unit sphere of $\CC^3$, so the amplitude vector $(c, b, a)$ of $\psi$ in the mode-ordered basis $(|100\rangle, |010\rangle, |001\rangle)$ is the third column of some $U \in SU(3)$. Since $G(U)$ fixes the vacuum and implements the mode rotation, $G(U)|001\rangle = G(U) f_3^\dagger |000\rangle = \sum_j U_{j3} f_j^\dagger |000\rangle = \psi$ exactly. (ii) The charge-sector distribution of any weight-1 state is $p = (0, 1, 0, 0)$, so the center piece $H(\{p_k\})$ vanishes; the quantum piece vanishes by Theorem~2. Hence $\SA = 0$ for $U(3)$, $SU(3)$, and $U(1)$ alike --- zero invariant entropy, whole curve, no sampling required. \hfill$\square$

\medskip

Numerical confirmation: all five stored curve points give $\max |\SA| = 1.9 \times 10^{-15}$, and the explicit $W = G(U)|001\rangle$ construction of Section~5.3 has residual $0.0$. The entire matching curve is gauge-equivalent to a single $\bar d$ quark in a rotated color frame, with invariant data identical to that of $|001\rangle$. \emph{The static coupling-matching analysis of \cite{companion}, Secs.~2--4, was matching gauge artifacts.} Ref.~\cite{companion} diagnosed its title question as ill-posed; the well-posed functional goes further and evaluates the would-be evidence at exactly zero.

\subsection{Haar Ensemble and Analytic Cross-Check}

Over $10^4$ Haar-random pure states: $S_{U(3)}$ has mean 1.5798, std 0.2391, range $[0.4103, 1.9984]$; $S_{SU(3)}$ has mean 1.3981, std 0.1617, range $[0.4066, 1.5849]$. These match the analytic values: for Haar states the sector probabilities are Dirichlet-distributed with parameters the sector dimensions, $p \sim \mathrm{Dir}(1, 3, 3, 1)$, giving $\mathbb{E}[S_{U(3)}] = 1.5767$ via digamma sums; merging the $SU(3)$-trivial weights $\{0, 3\}$ gives $\mathrm{Dir}(2, 3, 3)$ and $\mathbb{E}[S_{SU(3)}] = 1.3963$. Sampled-vs-analytic discrepancies $(0.003, 0.002)$ are consistent with $N = 10^4$ fluctuations.

\section{The Well-Posed Question and Its Answer}

\subsection{The Reformulation}

\begin{quote}
\textbf{Well-posed question.} Let $\Alg$ be the commutant of the gauge action on the one-generation ideal. Does there exist a \emph{natural} map from the single-copy gauge-invariant data $\EA(|\Omega\rangle\langle\Omega|)$ of an algebra-selected vacuum $|\Omega\rangle$ to the three Standard Model couplings $(\alpha_3, \alpha_2, \alpha_1)$ --- ``natural'' meaning the assignment of invariant parameters to gauge factors is dictated by the algebra, not chosen by hand?
\end{quote}

This is well-posed where the title question of \cite{companion} was not: $\EA$ is gauge-invariant by construction (verified to $10^{-15}$), and ``how many independent parameters are available'' is now a theorem rather than a frame choice. Under $U(3)$ the invariant data has exactly 3 real parameters --- so there is naive room, one per coupling. The answer is nevertheless no, for four independent reasons, each verified in the accompanying code.

\subsection{Reason 1: the Sectors Are Charge Sectors, Not Gauge Sectors}

The center of the $U(3)$-invariant algebra is spanned by the projectors onto eigenspaces of the \emph{single operator} $N$ (electric charge $Q = N/3$). The four labels $\{0, 1/3, 2/3, 1\}$ are values of one $U(1)$ quantum number. The $SU(3)$ factor contributes no additional center label --- weights 1 and 2 carry the $\mathbf{3}$ and $\bar{\mathbf{3}}$, but that is representation \emph{type}, not a continuous parameter --- and there are four sectors for three would-be couplings. A map $(p_1, p_2, p_3) \to (\alpha_3, \alpha_2, \alpha_1)$ requires a bijection $\{\text{charge sectors}\} \to \{\text{gauge factors}\}$ that nothing in the algebra provides: the invariant data simply does not factorize along gauge factors. This is the ``index ambiguity'' of \cite{companion} (Sec.~2.1 Remark there) promoted from labeling worry to structural fact.

\subsection{Reason 2: the Algebra-Canonical Vacuum Has Only Two Parameters}

The rotor vacuum's invariant data is exactly
\begin{equation*}
p = \left(\tfrac{1}{2}, \tfrac{3}{14}, \tfrac{3}{14}, \tfrac{1}{14}\right), \qquad \text{with } p_1 = p_2 \text{ exactly},
\end{equation*}
because the vacuum $(\sqrt{7}\,e_0 - iu)/\sqrt{14}$ has uniform $|\text{amplitude}|^2 = 1/14$ on all seven imaginary units, forcing $p_k$ proportional to the sector dimension on weights 1--3 (numerical check: $|p_1 - p_2| < 10^{-16}$; the rationals certified against the floats to $10^{-10}$). The canonical vacuum therefore carries \textbf{two independent invariant parameters --- fewer than three couplings.} The permutation degeneracy that killed every hierarchical reading in the frame-dependent analysis (\cite{companion}, Sec.~4.2 and Appendix~A) \emph{survives gauge-invariantization} as an exact sector degeneracy. The obstruction was never the frame; it is the algebra's symmetry. We stress that the parameter count is decisive jointly with the naturality requirement of Section~6.1 and the ansatz class inherited from \cite{companion} --- two invariant parameters cannot fill three coupling slots under any algebra-dictated assignment within that class --- not as bare dimension counting, which a sufficiently contrived map could always evade.

\subsection{Reason 3: No Assignment Reproduces the Hierarchy Anyway}

Even granting an unprincipled sector-to-coupling assignment, the numbers refuse. We tested every ordered assignment of three distinct charge sectors to the three couplings under three feature maps ($p_k$, $\log_2 p_k$, $-p_k \log_2 p_k$) --- $3 \times 24 = 72$ combinations, of which 60 yield a finite gap ratio (the $p_1 = p_2$ degeneracy renders 12 indeterminate) --- against the inherited two-parameter logarithmic map (\cite{companion}, Eq.~1), i.e.\ demanding $r = (x_a - x_b)/(x_b - x_c)$ equal $\rSM = 1.8174$ (\cite{companion}, Eqs.~2--3; GUT-normalized $\alpha_1$, a convention this target inherits). Result: \textbf{none lands within the widened tolerance 0.05 of \cite{companion}} (let alone the 0.01 of its Experiment 1); the best miss is
\begin{equation*}
r = 2 \text{ exactly} \quad (\text{feature } p,\ \text{descending sectors } (0, 1, 3) \text{ or } (0, 2, 3)), \qquad \text{distance } 0.1826.
\end{equation*}
We flag this near-miss explicitly so nobody rediscovers it as a ``signal'': $r = 2$ versus 1.8174 is a 10\% coincidence of small fractions --- $(1/2 - 3/14)/(3/14 - 1/14) = 2$ --- reachable only through assignments that, absurdly, put the \emph{neutrino} sector and the \emph{positron} sector on two different gauge couplings while skipping one of the two exactly degenerate quark sectors. Assignments using both quark sectors give $r \in \{0, -1, \pm\infty\}$, degenerate by $p_1 = p_2$. This is numerology bait, not structure. (The no-match conclusion is also normalization-independent: Appendix~C.3 reruns the identical 60-assignment zoo with the abelian coupling in hypercharge normalization, target $\rY = 1.0433$ instead of the GUT-normalized 1.8174, and again finds nothing within 0.05 --- with a \emph{worse} best miss, 0.2533.)

\subsection{Reason 4: No \texorpdfstring{$SU(2)_L$}{SU(2)L} Action Is Present in This One-Ideal Realization}

\textbf{Status: deliberately scoped lemma, verified in the accompanying code (Appendix~C.1), plus structural identification from Furey's construction.}

\medskip

\noindent\textbf{Lemma 4 (scoped).} \emph{Within this one-ideal, number-conserving three-mode realization, no $SU(2)_L$ action is present: the mode gauge group is exactly $U(3) \cong (SU(3)_c \times U(1)_{\mathrm{em}})/\mathbb{Z}_3$, and the centralizer of the color action inside its Lie algebra is the two-dimensional abelian $\mathrm{span}\{\mathbf{1}, N\}$ (Lemma~C.1). A two-ideal/chiral construction is required before the three-coupling question can even be posed.}

\medskip

The scope is deliberate, and we flag what the lemma does \emph{not} say: it is a statement about the present realization --- one generation ideal, number-conserving mode rotations --- not about Furey's framework, in which $SU(2)_L$ acts on the quaternionic factor of $\CC \otimes \mathbb{H} \otimes \mathbb{O}$, between ideals, on chiral doublets \cite{furey18}. Nothing here implies that $SU(2)_L$ cannot be represented division-algebraically; it implies that the arena in which the question of \cite{companion} was posed is too small to host it.

In detail: the mode-rotation (number-conserving Bogoliubov) gauge group of three Jordan--Wigner modes is exactly $U(3) \cong (SU(3)_c \times U(1)_{\mathrm{em}})/\mathbb{Z}_3$ --- there is no room for an independent $SU(2)$ acting on the color modes. This claim is not only a representation-theoretic argument but a lemma checked in the accompanying code: Appendix~C.1 verifies that the Hermitian generators preserving the mode span form a real Lie algebra of dimension exactly $10 = \dim u(3) + 1$, each exponentiating to $e^{i\varphi}G(U)$ with $U \in U(3)$ (residuals $< 3 \times 10^{-15}$); that the centralizer of the color action within this Lie algebra is the two-dimensional \emph{abelian} algebra $\mathrm{span}\{\mathbf{1}, N\}$, so no $su(2)$ commuting with $SU(3)_c$ exists; and that enlarging to number-violating Bogoliubov generators only adds directions that break the Hamming-weight (electric-charge) grading. More physically, the ideal $(\nu, \bar d_i, u_i, e^+)$ contains no weak-isospin doublet pair: $\nu_L$ pairs with $e^-_L$ (absent --- the ideal contains $e^+$), $u_L$ with $d_L$ (absent --- the ideal contains $\bar d$). Weak raising and lowering operators would map this ideal into a \emph{different} ideal; in Furey's full construction $SU(2)_L$ acts on the quaternionic factor of $\CC \otimes \mathbb{H} \otimes \mathbb{O}$ --- between ideals, never within one \cite{furey18}. And the toy's abelian charge is electromagnetic $Q = N/3$, not hypercharge: $SU(2)_L$ does not commute with $Q$, so any implementation would have to move states between Hamming-weight sectors, and the only weight-0/weight-3 pair, $(\nu, e^+)$, is not a doublet.

Consequently \emph{the three-coupling question cannot even be posed gauge-covariantly on a single generation ideal}: at most two gauge factors ($SU(3)_c$, $U(1)_{\mathrm{em}}$) act, and their joint invariant data is the charge distribution. Ref.~\cite{companion} asked whether three couplings are encoded in this system; the gauge-covariant accounting says the system never had three couplings' worth of symmetry to begin with.

\subsection{Verdict: More Constrained, Not Less}

The ill-posed version of the question was \emph{underconstrained}: a 14-real-parameter state space, a guaranteed-nonempty one-parameter matching manifold, statistically generic matches, W-class realizations (\cite{companion}, Secs.~2.2--2.3, 3.1, 4). The well-posed version is drastically \emph{more} constrained:

\begin{itemize}
\item the invariant state space shrinks from 14 parameters to 3 ($U(3)$);
\item the entire weighted-W matching mechanism is annihilated --- $\SA = 0$ on the whole curve;
\item the algebra-selected vacuum's data is fixed and degenerate: two parameters, $p_1 = p_2$;
\item no natural sector $\to$ factor map exists, and the exhaustive unnatural ones miss $\rSM$;
\item one of the three couplings has no corresponding symmetry in the toy at all.
\end{itemize}

\begin{quote}
\textbf{Answer to the re-posed question: no --- and now provably, not just diagnostically.} The single-copy gauge-invariant data of any pure state in this toy, under the full $U(3)$, is its electric-charge statistics. Charge statistics has no natural three-fold structure aligned with the three gauge factors; the canonical vacuum's statistics are two-parameter degenerate; and no $SU(2)_L$ action is present in this realization (Lemma~4). The coupling hierarchy is not, and cannot be, encoded in the gauge-invariant algebraic entropy --- nor, under $U(3)$, in any single-copy invariant data --- of a single pure state of this toy. What this closes and what it leaves open is delimited in Section~8: the entropy version and the single-copy version, not the full invariant-data version.
\end{quote}

\section{Discussion}

\subsection{The No-Go Ladder}

Between \cite{companion} and the present paper, the coupling-entropy question has now descended a complete ladder, each rung sharper than the last:

\begin{enumerate}
\item \textbf{Frame-dependent} (\cite{companion}, Secs.~2--4): the static matching analysis succeeded too easily --- matching was statistically generic, null controls matched equally well, and the one algebraically flavored success (the weighted-W curve) survived.
\item \textbf{Ill-posed} (\cite{companion}, Sec.~5.3): the quantity being matched, single-qubit entropy of the Jordan--Wigner factorization, is not gauge-invariant. Diagnosis, not yet answer.
\item \textbf{Well-posed} (this paper, Secs.~3--4): replace subsystem entropy by the entropy relative to the gauge-invariant algebra. The question now has a definite, frame-independent content: three real parameters under $U(3)$.
\item \textbf{Still no} (this paper, Sec.~6): the well-posed question has a negative answer for four independent structural reasons --- and it retroactively evaluates the rung-1 evidence at exactly zero (Sec.~5.4).
\end{enumerate}

The pattern is worth stating because it is the opposite of the usual fear about no-go results: making the question well-posed did not rescue a positive answer that ill-posedness had been obscuring; it converted a diagnostic ``the question is meaningless'' into a computed ``the question is meaningful and the answer is no.''

\subsection{All Edge, No Bulk: the Bridge to the Continuum --- With Its Dispute Included}

For pure states in this toy, the algebraic entropy is 100\% center term (Theorem~2) --- the discrete analog of the statement that a lattice-gauge entanglement entropy is carried by the classical boundary-flux (edge-mode) distribution. The lineage of that object in the continuum is well documented: Kabat's contact term in the conical entropy of gauge fields \cite{kabat}; the Donnelly--Wall edge-mode program, which interpreted it as the Shannon entropy of normal-flux sectors \cite{dw15, dw16}; Donnelly's lattice form, $S = H(\{p_R\}) + \sum_R p_R(\cdots)$, of which Eq.~(\ref{eq:salg}) is the abstract finite-dimensional version \cite{donnelly12, donnelly14}; and the Casini--Huerta--Rosabal algebra-with-center formulation that our definition discretizes directly \cite{chr}. The toy realizes this structure in its most extreme form: all edge, no bulk. Any future claim that vacuum entanglement encodes couplings in this framework is therefore, gauge-invariantly, a claim about precisely the object that literature studies as edge modes.

Honesty requires importing the dispute along with the structure. Casini, Huerta, Mag\'an, and Pontello \cite{chmp} argue that in the continuum the center/edge term is an artifact of the algebra (equivalently regularization) choice: local continuum algebras need not have centers, the classical term drops out of mutual information and relative entropy, and related work shows it is not distillable \cite{soni, moitra}. On that view, the quantity this paper computes wholesale --- $H(\{p_k\})$ --- is exactly the piece of ``entanglement entropy'' whose physical status is contested. Three things can be said in the toy's defense, none of them a refutation. First, in the toy the algebra is not a regularization choice: given the gauge group, the commutant is canonical --- the ambiguity that worries \cite{chmp} reappears here only as the (physically meaningful) choice of \emph{which group to gauge}, and we computed all three. Second, under a charge superselection rule the center term has a clean single-copy operational reading \cite{brs}: it is the Shannon entropy of the charge measurement that superselected observers can actually perform. Third, and cutting the other way: if one adopts the strict Casini--Huerta--Mag\'an--Pontello standard --- only mutual-information-like quantities are physical --- then the toy's gauge-invariant entropy is not merely classical but arguably \emph{empty}, and the no-go of Section~6 becomes if anything stronger. The dispute does not threaten our negative conclusion; it threatens only the consolation prize.

One further contrast pre-empts a natural question and frames the sequel. Donnelly's two-dimensional Yang--Mills results \cite{donnelly12, donnelly14} show that a gauge coupling \emph{is} exactly recoverable from the flux-sector (center) distribution --- the very object computed throughout this paper. But that recovery is \emph{reflection}, not determination: there, dynamics supplies a known one-parameter exponential family $p_R(g^2 A)$ of flux-sector probabilities, and a known family is statistically invertible given the theory --- the coupling can be read off the sector statistics because the theory says how it got in. The no-go of Section~6 concerns \emph{determination}: an algebra-dictated map from the invariant data of a single kinematic state to three coupling constants, with no dynamics supplying a family and one of the three gauge groups not acting on the construction at all. The two notions are complementary, not in tension --- sector statistics can perfectly well reflect a coupling that dynamics put there, even though no kinematic invariant data determines couplings here. The positive (reflection) case is developed exactly in \cite{ym2companion}, in two-dimensional Yang--Mills, where the coupling is not only reflected but efficiently estimable from the flux-sector statistics.

\subsection{Limitations}

The obvious one bears repeating: this is an 8-dimensional toy with no dynamics, no scale, and (Sec.~6.5) at most two gauge factors. The no-go results are exact but local to the construction; they constrain the minimal three-qubit literalism of the program of \cite{companion}, not gauge theory. The flow version of the question --- whether any RG-consistent trajectory on the sector simplex, with the rotor vacuum's $(1/2, 3/14, 3/14, 1/14)$ as boundary data, can reproduce one-loop running --- remains open, though Section~6.2 makes the expected outcome plain: the simplex has no natural three-coupling structure either. Finally, the $\rSM = 1.8174$ target inherits the GUT normalization $\alpha_1 = (5/3)\alpha_Y$ from \cite{companion}; the \emph{target} of the exhaustive search of Section~6.4 is in that sense convention-dependent, but the outcome was checked directly and is not: with the abelian coupling in hypercharge normalization, $\alpha_Y^{-1} = (5/3)\alpha_1^{-1} = 98.33$, the target becomes $\rY = 1.0433$, and the identical 60-assignment zoo again finds no match within 0.05 --- best miss 0.2533, farther than the GUT-normalized 0.1826 (Appendix~C.3). And the structural reasons (Secs.~6.2, 6.3, 6.5) never depended on the convention --- $p_1 = p_2$ does not care how the abelian coupling is normalized.

\section{Entropy Is Not the Full Invariant State}

This paper's no-go is, by design, a statement about gauge-invariant \emph{data}, and its headline quantity is an \emph{entropy}. This section delimits exactly how much each level closes, because the entropy is a lossy summary at two distinct stages --- and the second stage is the concrete next step for this program.

\textbf{Stage 1: the entropy is a lossy summary of the invariant state.} Corollary~2.1 is the sharpest possible illustration. The $U(1)$-invariant algebra (dimension 20) is five times larger than the $U(3)$-invariant algebra (dimension 4), and its reduced state $\EA(\rho)$ carries 11 real parameters against $U(3)$'s 3 --- yet for every pure state the two entropies are \emph{equal}, because the quantum piece vanishes identically (Theorem~2) and the two centers coincide. Vastly different invariant data, identical entropy. The moral generalizes: ``no coupling structure in the entropy'' does not by itself imply ``no coupling-like information in the invariant state.'' The entropy results of Section~5 close only the entropy version of the question; it is the reasons of Sections~6.2--6.5 --- which operate on the full single-copy data $\EA(\rho)$, not on its entropy --- that close the single-copy version.

\textbf{Stage 2: the single-copy state is not the full invariant ring.} The no-go of Section~6 governs single-copy gauge-invariant data --- everything measurable under the superselection rule \cite{brs}. The full polynomial invariant ring of a state under $G(U(3))$ is strictly larger, and for the rotor vacuum its low-degree elements are nonzero and exactly computable: with $v_1$, $v_2$ the weight-1 and weight-2 component vectors,
\begin{equation*}
|v_1 \wedge v_2| = \tfrac{1}{14} \text{ exactly}, \qquad
|\bar v_3 \cdot (v_1 \wedge v_2)| = 14^{-3/2} \approx 0.019090, \qquad
|\bar v_0\, v_3\, \overline{(v_1 \wedge v_2)}| = \tfrac{\sqrt{7}}{196} \approx 0.013499,
\end{equation*}
all verified gauge-invariant to $5.6 \times 10^{-17}$ under random $G(U)$. These are not functions of $(p_0, \ldots, p_3)$: they are constants of the gauge orbit invisible to $\EA(\rho)$, accessible only to joint measurements across multiple copies.

\textbf{What is closed, and what is not.} \emph{This paper closes the entropy version of the coupling question, and the single-copy invariant-data version --- not the full invariant-data version.} We state the residual scope plainly: no result here excludes coupling-like structure in the multi-copy invariant ring, beyond two observations --- the single-copy no-go does extend to everything operationally accessible under the superselection rule \cite{brs}, and we know of no principle that would route three coupling constants through multi-copy invariants while bypassing every single-copy observable.

\textbf{The concrete next step.} Characterizing the full $U(3)$ invariant ring on $\CC^8$ is finite, classical invariant theory, and undone: (i) compute the Hilbert series of the invariant ring (Molien--Weyl integration over $U(3)$, here a three-torus integral) to count independent invariants degree by degree; (ii) determine the number of functionally independent invariants directly, via the Jacobian rank of low-degree invariants at generic points; (iii) compare against the physically natural multi-copy observables --- Casimir moments and cumulants $\langle \psi^{\otimes m}|\, C \,|\psi^{\otimes m}\rangle$ for Casimirs $C$ of the diagonal $U(3)$ action on $(\CC^8)^{\otimes m}$, which organize the ring by copy number $m$. If a three-parameter coupling-like structure exists anywhere in this toy, it must live in (i)--(iii); Section~6 says only that it does not live in the single-copy data.

\section{The Two-Ideal Extension: Specification and Status}

The single most direct continuation --- and, by Lemma~4, the \emph{minimal} arena on which the three-coupling question is even posable --- is the two-ideal extension. We specify it here as a concrete computational program. Since this manuscript was first drafted, the computation has been completed; it is reported in a drafted companion manuscript \cite{twoideal}, where the obstruction-lifting construction and its gauge-invariant structure are computed. We do not import those results here beyond that one sentence: the specification below stands as the program was posed --- falsifiable as written, in either direction --- and the claims of the present paper are unchanged by it.

\begin{enumerate}

\item \textbf{Define the two-ideal Hilbert space.} Either the 16-dimensional doubling $\mathcal{H}_2 = \mathcal{H}_+ \oplus \mathcal{H}_-$ of the one-generation ideal by its conjugate ideal, or the full $\mathrm{Cl}(6)$ Fock construction carrying Furey's left/right ideal pair. In either case the arena must contain the chiral doublet pairs ($\nu_L$ with $e^-_L$, $u_L$ with $d_L$) whose absence from the single ideal is exactly what Lemma~4 formalizes.

\item \textbf{Specify the gauge action.} An explicit unitary representation $G_2$ of $SU(3)_c \times SU(2)_L \times U(1)_Y$ (or of the global form the construction actually delivers, e.g.\ $S(U(3) \times U(2))$-type) on $\mathcal{H}_2$, with $SU(2)_L$ implemented as the ideal-mixing quaternionic action and the abelian charge now \emph{hypercharge} rather than $Q$ --- verified to machine precision (homomorphism, mode-transformation property), exactly as in Section~2.

\item \textbf{Compute the commutant.} Run the identical machinery: SVD nullspace of the stacked commutator maps over Haar samples, certification on fresh samples, sandwich argument as in Lemma~C.2, and the analytic Schur decomposition of Appendix~B wherever the representation content can be identified. Output: the invariant algebra $\Alg_2$ and its Wedderburn structure.

\item \textbf{Decompose the center.} Identify the minimal central projections of $\Alg_2$ and the physical quantum numbers labeling them. The decisive structural question: do the sectors remain the level sets of a single abelian charge (the one-ideal outcome), or does a genuinely multi-factor sector lattice appear? Secondarily: do blocks with both $n_k, m_k > 1$ appear, so that pure states can carry a nonzero quantum piece --- genuine bulk, not only edge?

\item \textbf{Test for three-parameter structure.} Pose Section~6.1's well-posed question verbatim on $\Alg_2$: does the invariant data of an algebra-selected state factorize along the three gauge factors under an algebra-dictated (not hand-chosen) assignment? The one-ideal obstructions to beat are precise: the charge-sector reduction (Section~6.2) and the $p_1 = p_2$ degeneracy (Section~6.3).

\end{enumerate}

Either outcome of step 5 is informative: either the center acquires a natural threefold structure aligned with the gauge factors --- the first nontrivial opening for any version of the coupling--entropy correspondence in this program --- or the charge-sector obstruction persists, and the negative verdict extends to the smallest arena where the question is honestly posable. The companion manuscript \cite{twoideal} reports which.

\section{Conclusion}

Ref.~\cite{companion} ended by declaring the coupling-entropy question ill-posed until a gauge-invariant entanglement functional was specified, and named that specification the program's central open problem. This paper specified the functional --- the standard algebraic one, executed exactly --- and the completed program closes the question rather than reopening it. The well-posed question has a sharper negative answer than the ill-posed one: the gauge-invariant algebraic entropy of any pure state of the minimal fermionic toy is the Shannon entropy of its electric-charge statistics, entirely classical (``all edge, no bulk,'' Theorem~2); the algebra-canonical vacuum carries two independent invariant parameters, not three; the one positive-seeming static result of \cite{companion} is revealed to be a single quark in a rotated color frame, with invariant entropy exactly zero (Theorem~3); and no $SU(2)_L$ action is present in this one-ideal, number-conserving realization (Lemma~4), so the three-coupling question was never posable here in the first place.

For any version of the coupling-entropy correspondence within this division-algebraic program to survive, two things would need to exist. First, a richer arena: a two-ideal (or larger) construction on which $SU(2)_L$ acts, so that a three-coupling question is gauge-covariantly posable at all --- and on which one can then check whether the invariant center acquires a natural threefold structure aligned with gauge factors, or whether the charge-sector obstruction found here persists. Section~9 specifies that construction step by step; the computation has since been completed and is reported in a drafted companion manuscript \cite{twoideal}, whose results we do not import here. Second, a principle: an algebra-dictated, not hand-chosen, map from gauge-invariant sector statistics to coupling constants, presumably with scale dependence built in --- and no such principle currently exists. Absent it, the conclusion of the minimal case stands as computed: there is nothing in the gauge-invariant algebraic entropy of this toy --- nor in its single-copy invariant data --- for the Standard Model couplings to be encoded in; the multi-copy invariant ring of Section~8 is the one place in this toy where the question still lives. We offer the closed case, the machinery (which applies unchanged to any finite-dimensional gauge action), and the precise specification of what a successor construction must provide.

\appendix

\section{Reproducibility}

\textbf{Code.} All computations in this paper are produced by three files in the repository's \texttt{direction\_A/} directory:

\begin{itemize}
\item \texttt{alg\_entanglement.py} --- library: numerical commutant (SVD nullspace of the stacked commutator maps), Wedderburn structure identification (center, minimal central projections, per-block ranks), trace-preserving conditional expectation, algebraic entropy with the built-in cross-check $\SA = S(\EA(\rho)) - \sum_k p_k \log_2 m_k$ asserted to $10^{-8}$ on every call, Jordan--Wigner modes, the $G(U)$ construction, and Haar samplers.
\item \texttt{run\_direction\_A.py} --- every computation reported in Sections~4--6, in order; fixed seed 20260609; every claim asserted at runtime, so the script fails if any verification breaks; writes \texttt{results.json}.
\item \texttt{strengthen.py} --- the runtime-asserted lemmas of Appendix~C (the mode-gauge-group lemma with its Bogoliubov, particle-hole, and antilinear variants; the $SU(3)$-centralizer computation; the hypercharge-normalization rerun of the assignment zoo; and the cross-paper weighted-W point reconciliation of Section~5.2); fixed seed 20260609; writes \texttt{strengthen\_results.json}.
\end{itemize}

Run: \texttt{python run\_direction\_A.py} ($\approx$ 6 s; NumPy + SciPy only), then \texttt{python strengthen.py} for Appendix~C. The scripts read \texttt{../experiments/results/exp2\_results.json} (published with \cite{companion}) for the (five) stored weighted-W matching-curve amplitudes tested in Section~5.4.

\textbf{Sampling protocol.} Commutants: 50 Haar generator samples per gauge group, certified against 100 fresh samples (max commutator residual $< 1.2 \times 10^{-15}$) and checked for unital $*$-closure ($< 1.3 \times 10^{-15}$). Gauge invariance: 20 random gauge elements per state per group (assertion threshold $10^{-9}$; observed $< 3 \times 10^{-15}$). Haar ensemble: $10^4$ states, with a 200-state subsample run through the full block machinery (fast-formula agreement $2.0 \times 10^{-15}$; max $|$quantum piece$|$ $1.8 \times 10^{-15}$) and analytic Dirichlet cross-checks of the means asserted to 0.01.

\textbf{Exact rationals.} Sector probabilities reported as fractions ($1/2$, $3/14$, $3/14$, $1/14$, etc.) are produced by \texttt{Fraction(x).limit\_denominator($10^4$)} and asserted to agree with the floating-point values to $10^{-10}$; the closed-form entropies are asserted against the numerics to $10^{-12}$; $|v_1 \wedge v_2| = 1/14$ is asserted to $10^{-12}$. The statement $p_1 = p_2$ is exact by construction (uniform $|\text{amplitude}|^2 = 1/14$ on the seven imaginary units forces $p_k \propto$ sector dimension on weights 1--3) and checked numerically to $< 10^{-16}$.

\textbf{Construction residuals} (this run): CAR relations $0.0$; $G(U)$ unitarity $2.4 \times 10^{-15}$; homomorphism $6.8 \times 10^{-16}$; mode-transformation property $1.9 \times 10^{-15}$; $W = G(U)|001\rangle$ construction $0.0$, with $\EA$ reduction agreement $2.6 \times 10^{-16}$; $M_2$-block support identification asserted at threshold $10^{-9}$ --- an assertion threshold, not a measured residual (measured: $\leq 3.3 \times 10^{-15}$).

\section{Analytic Proof of Theorem 1}

The proof is three applications of Schur's lemma to the exterior-algebra decomposition of the Fock space. Throughout, $\omega = e^{2\pi i/3}$.

\subsection{The Fock Space Is the Exterior Algebra}

The map $f_{i_1}^\dagger \cdots f_{i_k}^\dagger |000\rangle \mapsto e_{i_1} \wedge \cdots \wedge e_{i_k}$ (for $i_1 < \cdots < i_k$, with $e_1, e_2, e_3$ the standard basis of $\CC^3$) is a unitary identification of the Hamming-weight-$k$ sector of $\CC^8$ with $\Lambda^k \CC^3$, and the canonical anticommutation relations make it consistent for arbitrary orderings: reordering creation operators produces exactly the sign of the corresponding wedge reordering. Under this identification, the mode-transformation property $G(U) f_b^\dagger G(U)^\dagger = \sum_a U_{ab} f_a^\dagger$ says precisely that $G(U)$ acts on $\Lambda^k \CC^3$ as the $k$-th exterior power $\Lambda^k U$. Hence, as a representation of $U(3)$,
\begin{equation}
\CC^8 \cong \Lambda^0 \CC^3 \oplus \Lambda^1 \CC^3 \oplus \Lambda^2 \CC^3 \oplus \Lambda^3 \CC^3, \qquad
G(U) \cong 1 \oplus U \oplus \Lambda^2 U \oplus \det U,
\label{eq:extalg}
\end{equation}
with sector dimensions $(1, 3, 3, 1)$ --- the blockwise definition of Section~2, now recognized as functorial rather than ad hoc.

\subsection{Schur's Lemma, in the Form Used}

Let a compact group act unitarily on $\mathcal{H} \cong \bigoplus_i V_i \otimes \CC^{m_i}$, where the $V_i$ are pairwise inequivalent irreducibles occurring with multiplicities $m_i$. Then the commutant of the action is $\bigoplus_i \mathbf{1}_{V_i} \otimes M_{m_i}(\CC)$: every intertwiner between inequivalent irreducibles vanishes, and every self-intertwiner of an irreducible is scalar. The commutant of each gauge action is therefore read off from two facts alone: which summands of Eq.~(\ref{eq:extalg}) are irreducible, and which are equivalent to one another.

\subsection{Gauging \texorpdfstring{$U(3)$}{U(3)}: Four Inequivalent Irreps, Multiplicity-Free}

Each $\Lambda^k \CC^3$ is irreducible under $U(3)$: $\Lambda^0$ and $\Lambda^3$ are one-dimensional; $\Lambda^1 = \CC^3$ is irreducible because $U(3)$ acts transitively on the unit sphere of $\CC^3$, so a nonzero invariant subspace contains every unit vector; and $\Lambda^2 \CC^3 \cong (\Lambda^1 \CC^3)^* \otimes \det$ as $U(3)$-representations (via the wedge pairing $\Lambda^2 \times \Lambda^1 \to \Lambda^3$), and the twist of the dual of an irreducible by a character is irreducible. The four are pairwise inequivalent already on the center of $U(3)$: the scalar $e^{i\theta}\mathbf{1}$ acts on $\Lambda^k$ as $e^{ik\theta}$, and the central characters $k = 0, 1, 2, 3$ are distinct. Eq.~(\ref{eq:extalg}) is thus multiplicity-free with four pairwise inequivalent summands, and Schur's lemma gives
\begin{equation*}
\{G(U) : U \in U(3)\}' = \CC P_0 \oplus \CC P_1 \oplus \CC P_2 \oplus \CC P_3 \cong \CC^4:
\end{equation*}
the four weight projectors and nothing else --- abelian, equal to its own center.

\subsection{Gauging \texorpdfstring{$SU(3)$}{SU(3)}: the Two Trivial Summands Merge}

Restricting to $\det U = 1$ collapses the central-character argument exactly where it should: $\Lambda^0$ and $\Lambda^3$ both become the trivial representation, which therefore occurs with multiplicity 2, on $\mathrm{span}\{|000\rangle, |111\rangle\}$. $\Lambda^1 = \mathbf{3}$ and $\Lambda^2 \cong \bar{\mathbf{3}}$ remain irreducible (the transitivity and duality arguments above use only $SU(3)$) and remain inequivalent: on the central element $\omega\mathbf{1} \in SU(3)$ their characters are $3\omega \neq 3\bar\omega$. Schur's lemma gives
\begin{equation*}
\{G(U) : U \in SU(3)\}' \cong M_2 \oplus \CC \oplus \CC \qquad (\dim = 6),
\end{equation*}
the $M_2$ block being the full matrix algebra on the multiplicity space $\mathrm{span}\{|000\rangle, |111\rangle\} = \mathrm{span}\{\nu, e^+\}$, with center $\CC^3$. That the $\nu$--$e^+$ coherence is gauge-invariant under color-only gauging is thus a one-line representation-theoretic fact: it is a coherence between two copies of the trivial representation.

\subsection{Gauging \texorpdfstring{$U(1)$}{U(1)}: Each Charge Sector Is Fully Degenerate}

$G(e^{i\theta}\mathbf{1}) = e^{i\theta N}$ acts on the weight-$k$ sector as the scalar $e^{ik\theta}$: each sector is the isotypic component of a one-dimensional character, occurring with multiplicity equal to the sector dimension $(1, 3, 3, 1)$, and the four characters are distinct. Schur's lemma gives everything block-diagonal in the weight grading,
\begin{equation*}
\{e^{i\theta N} : \theta \in \RR\}' = M_1 \oplus M_3 \oplus M_3 \oplus M_1 \qquad (\dim = 1 + 9 + 9 + 1 = 20),
\end{equation*}
with center $\CC^4$ --- the same four weight projectors as for $U(3)$, which is the representation-theoretic content of Corollary~2.1. \hfill$\square$

\subsection{Consistency with the Numerics}

{\sloppy The dimensions $(4, 6, 20)$, block supports (weights $\{0\}, \{1\}, \{2\}, \{3\}$; $\{0,3\}, \{1\}, \{2\}$; $\{0\}, \{1\}, \{2\}, \{3\}$), centers ($\CC^4$, $\CC^3$, $\CC^4$), and the location of the $M_2$ block on basis indices $\{0, 7\}$ agree item for item with the numerically computed and certified structures of Section~4 (commutator residuals $< 1.2 \times 10^{-15}$) and with the sandwich certification of Lemma~C.2. The numerics assumed none of the above --- the agreement is a genuine cross-check, not a restatement.\par}

\section{Code-Verified Lemmas}

Three statements used in the main text are recorded here as lemmas verified at runtime by the accompanying code, rather than as numerical observations or structural arguments. Code: \texttt{direction\_A/strengthen.py} (seed 20260609; NumPy + SciPy only); raw outputs in \texttt{direction\_A/strengthen\_results.json}; every claim below is asserted at runtime, so the script fails if any verification breaks.

\subsection{The Mode-Preserving Unitary Group Is \texorpdfstring{$\{e^{i\varphi} G(U) : U \in U(3)\}$}{\{exp(i phi) G(U) : U in U(3)\}}}

\textbf{Lemma C.1.} \emph{Let $F = \mathrm{span}_\CC\{f_1, f_2, f_3\} \subset M_8(\CC)$ be the span of the three Jordan--Wigner modes. Then:}

\emph{(i) (Lie algebra.) $L := \{h = h^\dagger : [h, f_k] \in F,\ k = 1, 2, 3\}$ has $\dim_\RR L = 10 = \dim u(3) + 1$, and $L = d\Gamma(u(3)) \oplus \RR\cdot\mathbf{1}$, where $d\Gamma(a) = \sum_{ij} a_{ij} f_i^\dagger f_j$. Every element of $L$ commutes with the number operator $N$.}

\emph{(ii) (Exponentiation.) For every $h \in L$, $W = e^{ih}$ satisfies $W f_k W^\dagger = \sum_j M_{kj} f_j$ with $M \in U(3)$, and $W = e^{i\varphi} G(M^\dagger)$ exactly.}

\emph{(iii) (Converse --- group level, all components.) Any unitary $W$ with $W F W^\dagger \subseteq F$ is of the form $e^{i\varphi} G(U)$: the canonical anticommutation relations force the induced coefficient matrix $M$ to be unitary, and $W G(M^\dagger)^\dagger$ commutes with every $f_k$ and $f_k^\dagger$, which generate all of $M_8(\CC)$, so it is scalar by Schur's lemma. No extra connected or disconnected components exist.}

\emph{(iv) (The $SU(2)_L$ obstruction proper.) The centralizer of the color action $\{G(V) : V \in SU(3)\}$ inside $L$ is $\mathrm{span}_\RR\{\mathbf{1}, N\}$: two-dimensional and abelian. An $SU(2)_L$ gauge factor would have to commute with $SU(3)_c$ and act on the modes; $su(2)$ needs a three-dimensional non-abelian algebra, and there is none.}

\textbf{Numerical checks.} The constraint $[h, f_k] \in F$ is linear in the 64 real parameters of a Hermitian $8 \times 8$ matrix; its solution space was computed as an SVD nullspace. The dimension count is unambiguous: largest ``zero'' singular value $1.1 \times 10^{-15}$ versus smallest nonzero singular value $1.0$. Constraint residual of the solution basis $9.3 \times 10^{-16}$; $\|[h, N]\| \leq 1.1 \times 10^{-15}$ for every solution; the identification $L = d\Gamma(u(3)) \oplus \RR\cdot\mathbf{1}$ is an exact rank check ($10 = 10$). Exponentiation was verified for all 10 basis elements and 20 random real combinations: $\max \|W f_k W^\dagger - \sum_j M_{kj} f_j\| = 2.8 \times 10^{-15}$, $\max \|MM^\dagger - \mathbf{1}\| = 4.7 \times 10^{-16}$, $\max \|W - e^{i\varphi}G(M^\dagger)\| = 2.5 \times 10^{-15}$ (all residuals $< 3 \times 10^{-15}$). The algebra generated by $\{f_k, f_k^\dagger\}$ has dimension $64 = \dim M_8(\CC)$, validating the Schur step. The $SU(3)$-centralizer in $L$ has dimension exactly 2, equals $\mathrm{span}\{\mathbf{1}, N\}$ by exact rank check, is abelian (basis commutator $< 10^{-12}$), and was certified on 20 fresh Haar $SU(3)$ samples (residual $1.1 \times 10^{-15}$). This grounds the scoped Lemma~4 of Section~6.5 in an explicit numerical check, not in a structural argument alone.

\textbf{Bogoliubov, particle-hole, and antilinear variants (also verified numerically).} Allowing number-violating maps $f_i \mapsto \sum_j (U_{ij} f_j + V_{ij} f_j^\dagger)$ --- constraint $[h, f_k] \in F \oplus F^\dagger$ --- enlarges the generator space to $\dim_\RR 16 = \dim so(6) + 1$, identified exactly as $d\Gamma(u(3)) \oplus \RR\cdot\mathbf{1}$ plus the six pair terms $f_i^\dagger f_j^\dagger + \mathrm{h.c.}$ and $i(f_i^\dagger f_j^\dagger - \mathrm{h.c.})$. But every added direction violates the Hamming-weight (electric-charge) grading: the number-conserving part collapses back to $L$ exactly (dimension 10, identical span); the pair generator $f_1^\dagger f_2^\dagger + \mathrm{h.c.}$ has $\|[h, N]\| = 4.0$, its commutator with $f_1$ has a component of norm 2.0 entirely outside $F$, and its exponential has off-weight-block norm 0.84 (it moves states across charge sectors). A particle-hole unitary exists --- $X \otimes X \otimes X$ implements $f_k \mapsto (f_1^\dagger, -f_2^\dagger, f_3^\dagger)$ exactly (residual 0) --- but it inverts the charge grading ($N \mapsto 3 - N$) and swaps $F \leftrightarrow F^\dagger$, so it lies outside the span-preserving group; $SU(2)$, being connected, cannot live in a disconnected coset. Antilinear maps add nothing new: all $f_k$ are real matrices, so every antiunitary preserving $F$ is entrywise conjugation $K$ composed with a span-preserving unitary, and $K G(U) K = G(\bar U)$ (residual 0).

\subsection{The Commutant Identifications Are Exact (Sandwich Argument)}

\textbf{Lemma C.2.} \emph{For each gauge choice $\mathcal{G} \in \{U(3), SU(3), U(1)\}$, the commutant $\{G(g) : g \in \mathcal{G}\}'$ is exactly the algebra identified in Section~4: the abelian $\CC^4$ of weight projectors (dimension 4), $M_2 \oplus \CC \oplus \CC$ with the $M_2$ block on $\mathrm{span}\{|000\rangle, |111\rangle\}$ (dimension 6), and $M_1 \oplus M_3 \oplus M_3 \oplus M_1$ of all weight-preserving operators (dimension 20), respectively.}

\textbf{Proof (sandwich).} For each gauge choice, exhibit elements that commute with every $G(g)$ \emph{exactly, by construction}: the four weight projectors ($U(3)$); the weight-1 and weight-2 projectors together with the four matrix units on $\mathrm{span}\{|000\rangle, |111\rangle\}$, both weights being $SU(3)$-trivial ($SU(3)$); and all weight-preserving matrix units ($U(1)$). These span spaces of dimensions 4, 6, and 20, giving lower bounds on the commutant dimensions. The numerically computed commutant --- the SVD nullspace of the stacked commutator maps over the Haar samples --- \emph{contains} the true commutant, since commuting with the whole group implies commuting with the samples; its dimension is therefore an upper bound. The computed dimensions are 4, 6, and 20: the bounds coincide, forcing equality. The commutant is exactly the span of the exhibited elements, and the numerical identification is exact rather than approximate. \hfill$\square$

All three identifications also follow from Schur's lemma alone --- that is the analytic proof of Theorem~1 in Appendix~B. Lemma~C.2 is logically independent of it: the sandwich uses only exhibited commuting elements and the numerically computed dimension, so the two proofs certify each other.

\subsection{The Zoo Miss Is Normalization-Independent}

\textbf{Lemma C.3 (computational).} \emph{With the abelian coupling in hypercharge normalization, $\alpha_Y^{-1} = (5/3)\alpha_1^{-1} = 98.33$, the target gap ratio becomes $\rY = 1.0433$ in place of the GUT-normalized $\rSM = 1.8174$ --- and the identical 60-assignment sector-to-coupling zoo of Section~6.4, run on the rotor vacuum's exact invariant data $p = (1/2, 3/14, 3/14, 1/14)$ against $\rY$ with the same 0.05 window, again finds nothing: no assignment within 0.05, none within 0.01. The conclusion of Section~6.4 is normalization-independent.}

\textbf{Numerical checks.} Inputs as in \cite{companion}, from \cite{pdg}: $\alpha_3 = 0.1179$, $\alpha_2 = 0.03374$, $\alpha_1 = 0.01695$ ($\alpha_1^{-1} = 58.9971$); $\alpha_Y^{-1} = (5/3) \times 58.9971 = 98.3284$, giving $\rY = 1.043286$ (rounded-table variant 98.3333, $\rY = 1.043508$). The reimplemented zoo reproduces the published GUT-target best miss (0.18256) as a cross-check. Against $\rY$, the best miss is 0.2533, at $r = \log_2 3 / \log_2(7/3) = 1.296607$ (feature $\log_2 p_k$, sectors $(e^+, \bar d, \nu)$ and its exact mirror $(e^+, u, \nu)$, degenerate by $p_1 = p_2$); rounded-target variant 0.2531. The $\alpha_Y$ miss is \emph{worse} than the GUT miss (0.1826). Assignments using both quark sectors still give $r \in \{0, -1, \pm\infty\}$: the $p_1 = p_2$ degeneracy and the four-sectors-versus-three-couplings mismatch are untouched by the choice of abelian normalization, as expected.

\begin{thebibliography}{20}

\bibitem{companion} W.T. Trevena, ``No-Go Results for a Three-Qubit Entropy Ansatz for Gauge-Coupling Hierarchies, with an Exact Octonionic Rotor Construction in an Appendix,'' companion paper (2026); code and manuscript at the same repository.

\bibitem{furey14} C. Furey, ``Generations: three prints, in colour,'' JHEP 2014, 046 (2014).

\bibitem{furey18} C. Furey, ``$SU(3)_C \times SU(2)_L \times U(1)_Y$ ($\times U(1)_X$) as a symmetry of division algebraic ladder operators,'' Eur. Phys. J. C 78, 375 (2018).

\bibitem{stoica} O.C. Stoica, ``Leptons, quarks, and gauge from the complex Clifford algebra $\mathrm{Cl}(6)$,'' Adv. Appl. Clifford Algebras 28, 52 (2018); arXiv:1702.04336.

\bibitem{szangolies} J. Szangolies, ``The Standard Model Symmetry and Qubit Entanglement,'' Entropy 27(6), 569 (2025); arXiv:2512.17328.

\bibitem{jordanwigner} P. Jordan and E. Wigner, ``\"Uber das Paulische \"Aquivalenzverbot,'' Z. Phys. 47, 631 (1928).

\bibitem{zanardi} P. Zanardi, ``Virtual quantum subsystems,'' Phys. Rev. Lett. 87, 077901 (2001).

\bibitem{barnum} H. Barnum, E. Knill, G. Ortiz, and L. Viola, ``A subsystem-independent generalization of entanglement,'' Phys. Rev. Lett. 92, 107902 (2004).

\bibitem{chr} H. Casini, M. Huerta, and J.A. Rosabal, ``Remarks on entanglement entropy for gauge fields,'' Phys. Rev. D 89, 085012 (2014); arXiv:1312.1183.

\bibitem{donnelly12} W. Donnelly, ``Decomposition of entanglement entropy in lattice gauge theory,'' Phys. Rev. D 85, 085004 (2012); arXiv:1109.0036.

\bibitem{donnelly14} W. Donnelly, ``Entanglement entropy and nonabelian gauge symmetry,'' Class. Quantum Grav. 31, 214003 (2014); arXiv:1406.7304.

\bibitem{kabat} D. Kabat, ``Black hole entropy and entropy of entanglement,'' Nucl. Phys. B 453, 281 (1995); arXiv:hep-th/9503016.

\bibitem{dw15} W. Donnelly and A.C. Wall, ``Entanglement entropy of electromagnetic edge modes,'' Phys. Rev. Lett. 114, 111603 (2015); arXiv:1412.1895.

\bibitem{dw16} W. Donnelly and A.C. Wall, ``Geometric entropy and edge modes of the electromagnetic field,'' Phys. Rev. D 94, 104053 (2016); arXiv:1506.05792.

\bibitem{chmp} H. Casini, M. Huerta, J.M. Mag\'an, and D. Pontello, ``Logarithmic coefficient of the entanglement entropy of a Maxwell field,'' Phys. Rev. D 101, 065020 (2020); arXiv:1911.00529.

\bibitem{brs} S.D. Bartlett, T. Rudolph, and R.W. Spekkens, ``Reference frames, superselection rules, and quantum information,'' Rev. Mod. Phys. 79, 555 (2007).

\bibitem{soni} R.M. Soni and S.P. Trivedi, ``Aspects of Entanglement Entropy for Gauge Theories,'' JHEP 01 (2016) 136; arXiv:1510.07455 (non-distillability of the classical edge term). See also S. Ghosh, R.M. Soni, and S.P. Trivedi, ``On the Entanglement Entropy for Gauge Theories,'' JHEP 09 (2015) 069, arXiv:1501.02593, and, for the continuum (3+1)-d free $U(1)$ case, R.M. Soni and S.P. Trivedi, ``Entanglement Entropy in (3+1)-d Free U(1) Gauge Theory,'' JHEP 02 (2017) 101, arXiv:1608.00353.

\bibitem{moitra} U. Moitra, R.M. Soni, and S.P. Trivedi, ``Entanglement Entropy, Relative Entropy and Duality,'' JHEP 08 (2019) 059; arXiv:1811.06986 (the classical term does not contribute to relative entropy or mutual information in the continuum limit).

\bibitem{pdg} Particle Data Group, ``Review of Particle Physics,'' Prog. Theor. Exp. Phys. 2022, 083C01 (2022) --- coupling values entering $\rSM = 1.8174$ via \cite{companion}, Eqs.~(2)--(3).

\bibitem{ym2companion} W.T. Trevena, ``Fisher Information and Coupling Reconstruction from Edge-Sector Statistics in Two-Dimensional Yang--Mills,'' companion paper (2026); code and manuscript at the same repository.

\bibitem{twoideal} W.T. Trevena, ``Room for Three Couplings, and Only the Room: The Gauge-Invariant Algebra of a One-Generation Chiral Fock Space,'' companion manuscript (2026); code and manuscript at the same repository.

\end{thebibliography}

\end{document}
